\scanpage{117}
\uline{Notations}\quad Pour \(X\) variété projective non singulière.
\[
\left\{\begin{array}{ll}
H^*(X).&\text{cohomologie }\ell\text{-adique}\\
P^*(X)&\text{Partie primitive de la cohomologie}
\end{array}\right.
\]
\margin{(on peut mettre des coeff.}
\[
\left\{\begin{array}{ll}
A^*(X)=A_r^*(X)&\text{anneau de Chow (équiv. rationnelle)}\\
N^*(X)=A_a^*(X)&\text{``\quad pour équiv. algébrique}\\
J^*(X)=A_r^*(X)^a&
\end{array}\right.
\]
\diagram{\begin{array}{ccccccccc}0&\longrightarrow&J^*(X)&\unsure{\text{---}|\longrightarrow}&A^*(X)&\longrightarrow&N^*(X)&\unsure{\text{---}|\longrightarrow}&0\end{array}}
\[
\left\{\begin{array}{ll}
M^*(X)&\text{partie primitive de }N^*(X)\\
I^*(X)&\text{partie primitive de }J^*(X)
\end{array}\right.
\]
\uline{N.B.} Il \unsure{ne|me} semble pas qu'il y ait une définition raisonnable de la ``partie primitive'' de \(A^*(X)\) \(\text{---}\)
\margin{116}
